\documentclass[11pt]{article}
\usepackage[T1]{fontenc}
\usepackage{lmodern}
\usepackage[margin=1in]{geometry}
\usepackage{amsmath,amssymb,amsthm,mathtools}
\usepackage{microtype}
\usepackage{xcolor}
\PassOptionsToPackage{hyphens}{url}
\usepackage{hyperref}
\usepackage{bookmark}
\hypersetup{colorlinks=true,linkcolor=blue!45!black,citecolor=blue!45!black,
  urlcolor=blue!45!black,pdfauthor={OpenAI},
  pdftitle={A Counterexample to the Group-Ring Determinant Conjecture}}
\input{glyphtounicode-cmex.tex}
\pdfgentounicode=1
\newtheorem{theorem}{Theorem}[section]
\newtheorem{proposition}[theorem]{Proposition}
\newtheorem{lemma}[theorem]{Lemma}
\theoremstyle{remark}
\newtheorem{remark}[theorem]{Remark}
\numberwithin{equation}{section}
\newcommand{\F}{\mathbb F}
\newcommand{\Z}{\mathbb Z}
\newcommand{\Q}{\mathbb Q}
\newcommand{\R}{\mathbb R}
\newcommand{\C}{\mathbb C}
\newcommand{\Mat}{\operatorname{M}}
\newcommand{\diag}{\operatorname{diag}}
\newcommand{\N}{\mathcal N}
\setlength{\emergencystretch}{2em}
\title{A Counterexample to the Group-Ring Determinant Conjecture}
\author{OpenAI}
\date{September 23, 2026}

\begin{document}
\maketitle
\begin{abstract}
We disprove the unrestricted group-ring Determinant Conjecture. We construct a finitely generated group $G$ and a square matrix over $\mathbb Z[G]$ that is invertible over $\mathbb Q[G]$ and has Fuglede--Kadison determinant strictly between zero and one. The logarithmic integral defining the determinant is finite.
\end{abstract}

\section{Introduction and proof overview}
\label{sec:introduction}

A nonsingular integer matrix has determinant of absolute
value at least one. The group-ring Determinant Conjecture asks for an
analogue when the entries are finite integer combinations of elements
of a discrete group. The relevant determinant, introduced by Fuglede
and Kadison~\cite{FK1952}, uses the trace of a logarithm rather than
a finite-dimensional determinant polynomial.

Here are the conventions that fix this question. Let $G$ be a discrete
group and let $\Z[G]$ denote its ordinary group ring of finite sums.
Use the left regular representation
\[
 (L_g\xi)(h)=\xi(g^{-1}h),\qquad \xi\in\ell^2(G).
\]
For $A\in\Mat_n(\Z[G])$, write $T_A$ for its operator on
$\ell^2(G)^n$ and $e_i$ for the vector with $\delta_1$ in component
$i$. The probability spectral measure of the positive square is
\[
 \mu_A(E)=\frac1n\sum_{i=1}^n
 \langle \mathbf1_E(T_A^*T_A)e_i,e_i\rangle
 \quad(E\subseteq[0,\infty)\text{ Borel}).
\]
The square-matrix instance of the conjecture is
\begin{equation}\label{eq:conjecture}
 \int_{(0,\infty)}\log t\,d\mu_A(t)\ge0.
\end{equation}
The zero atom is omitted and the general integral may be $-\infty$.
The usual conjecture also permits rectangular matrices
\cite[Conjecture 13.2]{Luck2002}; a square counterexample already
contradicts it. Our example will have a bounded inverse, so its
spectrum stays away from zero.

Write $\N(G)$ for the group von Neumann algebra and use the
\emph{unnormalized} matrix trace
$\tau_n(T)=\sum_i\langle Te_i,e_i\rangle$, so $\tau_n(I_n)=n$.
For invertible $T_A$ the determinant convention is
\begin{equation}\label{eq:normalization}
 \log\det_{\N(G)}(T_A)
 =\frac12\tau_n\bigl(\log(T_A^*T_A)\bigr)
 =\frac n2\int_{(0,\infty)}\log t\,d\mu_A(t).
\end{equation}
This extends the absolute value of the ordinary determinant.
The inversion unitary $\delta_g\mapsto\delta_{g^{-1}}$ identifies
our integral-matrix convention with right multiplication by $A^*$
on row vectors. It fixes the identity vectors, and therefore
preserves the corresponding positive-square spectral measures.

\begin{theorem}\label{thm:main}
There are a finitely generated discrete group $G_D$, an integer
$n\ge1$, and $A\in\Mat_n(\Z[G_D])$ invertible over $\Q[G_D]$ such
that
\[
 0<\det_{\N(G_D)}(T_A)<1.
\]
The integral in Equation~\eqref{eq:conjecture} is finite and strictly
negative. Thus the unrestricted group-ring Determinant Conjecture
has a negative resolution.
\end{theorem}

\paragraph{The exact existence input.}
A unital ring is directly finite if $ab=1$ implies $ba=1$.
We use the following scalar failure of direct finiteness from the
complete characteristic-two companion.

\begin{theorem}[{\cite[Theorem 1.1]{CharacteristicTwo}}]
\label{thm:input}
There exist a finite field $K_2$ of characteristic two, a finitely
generated group $H_2$, and finite sums $a_2,b_2\in K_2[H_2]$ with
$a_2b_2=1$ and $b_2a_2\ne1$.
\end{theorem}

The companion's Section~6 specifies its field, group presentation
and scalar sums by a terminating finite prescription. We use the
existence theorem, not an executed numerical instance or a
group-word decision procedure. The coefficient field is finite but
is not asserted to be $\F_2$. There is no torsion-free hypothesis;
the companion construction uses torsion. All steps after this exact
input are proved below. In particular, a result only in odd
characteristic would not supply the binary premise used here.

\paragraph{Context and mathematical ancestry.}
Determinant bounds are part of the theory of $L^2$-invariants
\cite{Luck2002}. Approximation and closure arguments established
such bounds for large classes of groups, including amenable groups
\cite[Theorem 1.21, with its erratum]{Schick2001}; Elek and Szab\'o proved the
determinant bound for sofic groups
\cite[Theorem 5]{ElekSzabo2005}. Their proof uses finite graph
approximations and the integrality of products of nonzero eigenvalues
of positive integral matrices. Our construction instead starts from
failure of direct finiteness in characteristic two. The cited sofic theorem implies, as a
separate consequence, that its resulting group is not sofic.

The lower-bound question is distinct from determinant
\emph{approximation}. Kammeyer's counterexample to a finite-quotient
limiting formula has determinant $2$ over the discrete Heisenberg
group, while its principal congruence quotients of levels $3^i$, $i\ge1$,
have determinants at most $5^{1/3}$
\cite[Theorem 2 and Section 2]{Kammeyer2026}. That discrepancy does not violate
the lower bound one.

The finite-coordinate projection construction below is related to
the commuting wreath-product projections of Dicks and Schick
\cite[Section 3]{DicksSchick2002} and of Lehner, Neuhauser and
Woess~\cite[Section 2]{LehnerNeuhauserWoess2008}. Here the
projections separate one coefficient of a binary defect; we do not
import a spectral-distribution theorem from those works. The
analytic step is the classical tracial logarithmic derivative
\cite[Lemma 2]{FK1952}, also present in the path-integral determinant
viewpoint of de la Harpe and Skandalis
\cite[Lemma 1 and Proposition 2(d)]{deLaHarpeSkandalis1984}. We give its proof and its
application to block shears locally.

\paragraph{How a binary defect produces a real trace gap.}
Set $d=[K_2:\F_2]$. Representing multiplication by field elements
on the $d$-dimensional binary vector space $K_2$ gives matrices
$X,Y\in\Mat_d(\F_2[H_2])$ with $XY=I_d$ and a nonzero idempotent
$C=I_d-YX$. One entry of $C$ has a nonzero coefficient. We enlarge
the group to
\[
 G_D=\left(\bigoplus_{h\in H_2}C_3\right)\rtimes H_2,
\]
where $H_2$ acts by left translation on the coordinate labels. Averaging over the
cyclic group at one coordinate and taking complementary averages
at finitely many other coordinates isolates that coefficient.
The result is a rational orthogonal projection $P$ of size $d+1$
with trace $d+\theta$, for a positive rational number $\theta$.
Its reduction modulo two nevertheless factors through $d$ columns.
Section~\ref{sec:projection} proves the coefficient separation and
this factorization in full.

\paragraph{How the gap becomes an integral determinant.}
For a general rational projection $P$ of size $m$, suppose an odd
integer $D$ clears its denominators and $\bar P=uv$ factors through
$r$ columns over $\F_2[G]$. Choose a power $q$ of two congruent to
one modulo $D$. Section~\ref{sec:lifting} lifts the factorization modulo $q$
by a finite geometric series in a corner ring. Arbitrary integral
lifts then give an integral block matrix with an explicit rational
inverse. Section~\ref{sec:analysis} removes its two block shears
without changing the trace of the logarithm of its positive square.
The remaining diagonal matrix expands the first $r$ coordinate
summands by $q$ and contracts the range of
$P$ by $q^{-1}$, giving
\[
 \log\det(T_A)=(r-\tau_m(P))\log q.
\]
Finally Section~\ref{sec:assembly} uses $r=d$ and
$\tau_{d+1}(P)=d+\theta$ to obtain the exact value $q^{-\theta}$.
The projection construction, modular lifting and trace calculation
have distinct roles; none replaces the scalar existence input.

\section{From the scalar defect to a projection}
\label{sec:projection}

We first produce the trace gap. All group-ring sums in this section
are finite, and matrix multiplication keeps the indicated order.
For a group-ring sum $c$, write $[1]c$ for its identity coefficient.
The involution $*$ inverts group elements and conjugates coefficients;
on matrices it also transposes the indices. Thus an orthogonal
projection satisfies $P=P^*=P^2$ in the regular representation.

\begin{lemma}\label{lem:binary}
Let $K_2$ be a finite field of characteristic two and $H_2$ a group.
For $d=[K_2:\F_2]$, there is an injective unital ring homomorphism
\[
 \rho:K_2[H_2]\longrightarrow\Mat_d(\F_2[H_2]).
\]
For $a_2,b_2\in K_2[H_2]$ with $a_2b_2=1\ne b_2a_2$, their images $X,Y$ satisfy
$XY=I_d$ and $C:=I_d-YX\ne0$, with
\begin{equation}\label{eq:defect}
 C^2=C,\qquad XC=0,\qquad CY=0.
\end{equation}
\end{lemma}

\begin{proof}
Choose a binary basis of $K_2$. The matrix of multiplication by
$\lambda\in K_2$ is a unital ring representation; it is injective
because its value on $1\in K_2$ recovers $\lambda$. Apply this map
coefficientwise to a group-ring sum. The coefficient matrices
commute with the formal group symbols, so products are preserved.
Comparing the coefficient matrix at each group element proves
injectivity. Thus $XY=I_d$ and $I_d-YX\ne0$. Now
$(YX)^2=Y(XY)X=YX$, $X(I_d-YX)=X-(XY)X=0$ and
$(I_d-YX)Y=Y-Y(XY)=0$, proving all the identities.
\end{proof}

\begin{proposition}\label{prop:separation}
Let $X,Y\in\Mat_d(\F_2[H_2])$, $d\ge1$, satisfy
$XY=I_d\ne YX$. In the group
$G_D=(\bigoplus_{h\in H_2}C_3)\rtimes H_2$, where $H_2$ acts
by left translation on coordinate labels, there is an orthogonal
projection $P\in\Mat_{d+1}(\Z[1/3][G_D])$ and a finite set
$S\subseteq H_2\setminus\{1\}$ such that
\[
 \tau_{d+1}(P)=d+\theta,\qquad
 \theta=\frac13\left(\frac23\right)^{|S|}>0,
 \qquad \bar P=uv,
\]
where $u$ has size $(d+1)\times d$ and $v$ has size $d\times(d+1)$
over $\F_2[G_D]$. The integer $D=3^{1+|S|}$ clears all coefficients
of $P$.
\end{proposition}

\begin{proof}
Put $C=I_d-YX$. Choose an entry $C_{ij}$ and a group element
$h_0$ whose coefficient in that entry is nonzero. The coefficient
is $1$ because the field is $\F_2$. Left multiplication by
$h_0^{-1}$ therefore gives
\begin{equation}\label{eq:support}
 h_0^{-1}C_{ij}=1+\sum_{s\in S}s
\end{equation}
for a finite $S\subseteq H_2\setminus\{1\}$, possibly empty.

To isolate the identity term in Equation~\eqref{eq:support}, we
seek a nonzero rational projection $p$ with $psp=0$ for every
$s\in S$. The cyclic coordinates provide complementary projections
that enforce these vanishing products. Let $z_h$ generate the cyclic
coordinate labelled by $h$ and put
\[
 p_h=\frac{1+z_h+z_h^2}{3},\qquad
 p=p_1\prod_{s\in S}(1-p_s).
\]
The $p_h$ are self-adjoint idempotents: averaging the nine products
in $(1+z_h+z_h^2)^2$ gives three copies of each element of $C_3$.
They commute because distinct coordinates commute. Consequently
$p=p^*=p^2$, and the displayed product shows that $3^{1+|S|}p$
is integral. This is the elementary commuting-coordinate method
in the works cited above; compare
\cite[Lemmas 3.1 and 3.4]{DicksSchick2002}.

Conjugation obeys $s p_hs^{-1}=p_{sh}$. For $s\in S$, the
product $sps^{-1}$ contains $p_s$, while $p$ contains $1-p_s$.
All these coordinate factors commute, so
\begin{equation}\label{eq:separation}
 psp=p(sps^{-1})s=0.
\end{equation}
This argument does not need disjointness of the two collections of
coordinate labels. Inverse pairs or finite-order elements in $S$
cause no change.

The original coordinates $\{1\}\cup S$ are distinct. The coefficient
of the identity in their product is therefore the product of the
individual identity coefficients. Hence
\begin{equation}\label{eq:theta}
 \tau_1(p)=[1]p=\frac13\prod_{s\in S}\frac23=\theta>0.
\end{equation}
For empty $S$ this means $p=p_1$ and $\theta=1/3$.

The map $\Z[1/3]\to\F_2$ is defined because $3$ is odd. Reduce
the projection modulo two and combine
Equations~\eqref{eq:support} and~\eqref{eq:separation} to obtain
\[
 \bar p h_0^{-1} C_{ij}\bar p=\bar p.
\]
Let $f$ be the row of length $d$ with entry $\bar p h_0^{-1}$
at position $i$ and zero elsewhere. Let $w$ be the column with
entry $\bar p$ at position $j$ and zero elsewhere. Thus
$fCw=\bar p$. Define
\[
 P=\diag(I_d,p),\qquad
 u=\begin{pmatrix}X\\fC\end{pmatrix},\qquad
 v=\begin{pmatrix}Y&Cw\end{pmatrix}.
\]
Using Equation~\eqref{eq:defect} in this order gives
\[
 uv=\begin{pmatrix}XY&XCw\\fCY&fC^2w\end{pmatrix}
   =\diag(I_d,\bar p)=\bar P.
\]
The rational trace is $d+\theta$ by Equation~\eqref{eq:theta}.
It is the rational orthogonal projection, not its binary reduction,
that has this strictly positive extra trace.
\end{proof}

We have obtained a projection of real trace greater than $d$ but a
binary factorization through $d$ columns. The next step produces
a factorization through the same number of columns modulo a power of two, so
that it can be inserted into an integral matrix.

\section{Lifting the factorization and preserving integrality}
\label{sec:lifting}

We now work with any group $G$, an orthogonal rational projection
$P\in\Mat_m(\Q[G])$, and a factorization of $\bar P$ through $r$
columns over $\F_2[G]$. Assume that an odd integer $D$ clears the
denominators of $P$. Here orthogonality means $P^2=P=P^*$, where
$*$ transposes the matrix and inverts group elements.

For a power $q>1$ of two, the trace gap suggests the absolute value
of the operator we want: expand $r$ coordinate summands by $q$ and
contract the range of $P$ by $q^{-1}$. Its
logarithmic determinant should be $(r-\tau_m(P))\log q$, which is
negative when $\tau_m(P)>r$. The difficulty is to realize this
determinant by an integral matrix. We will add block shears whose
off-diagonal entries come from a factorization of $P$ modulo $q$.
The analytic justification that these shears preserve the determinant
is given in Section~\ref{sec:analysis}.

The first lemma lifts the binary factorization to every
power $q$ of two; only idempotence is needed for this step.

\begin{lemma}\label{lem:lift}
Let $G$ be a group, $m,r\ge1$, $D$ an odd positive integer, and
$P\in\Mat_m(\Q[G])$ satisfy $P^2=P$ and
$DP\in\Mat_m(\Z[G])$. Suppose $\bar P=uv$ over $\F_2[G]$, where $u$ and $v$
have respective sizes $m\times r$ and $r\times m$. For $q=2^k$,
$k\ge1$, the reduction $e=P_q$ has a factorization
$e=u_qv_q$ of the same sizes over $(\Z/q\Z)[G]$.
\end{lemma}

\begin{proof}
Since $D$ is odd, every coefficient of $P$ has a well-defined
reduction modulo $q$. Set $\mathcal R=(\Z/q\Z)[G]$ and choose
arbitrary finite lifts $u_0,v_0$ of $u,v$ over $\mathcal R$.
In the corner ring $e\Mat_m(\mathcal R)e$, whose identity is $e$,
put
\[
 Z=eu_0v_0e,\qquad W=e-Z.
\]
Modulo two, $Z$ is $\bar Puv\bar P=\bar P^3=\bar P$.
Every coefficient of $W$ is therefore even. A product of $k$
such coefficients is divisible by $2^k$, giving $W^k=0$.
This is nilpotence in the even ideal of the corner ring; the
corner ring itself need not be nilpotent.

Because $Z=e-W$, the finite geometric sum
\[
 Q=e+W+\cdots+W^{k-1}
\]
satisfies $QZ=ZQ=e-W^k=e$. When $k=1$, this sum is just $e$.
Thus $Q$ is a corner inverse of $Z$, not a claimed inverse in
the entire matrix ring. Set
\[
 u_q=Qeu_0,\qquad v_q=v_0e.
\]
Then $u_qv_q=Qeu_0v_0e=QZ=e$. Only finite sums and products
were used, so the resulting supports remain finite.
\end{proof}

For the integral construction, choose $q=2^k>1$ with
$q\equiv1\pmod D$. Such a power exists because two is a unit
modulo the odd integer $D$; for $D=1$, any $2^k>1$ works.
Choose integral lifts $U,V$ of the factors from Lemma~\ref{lem:lift}.
The diagonal operator with the desired expansion and contraction is
\[
 J=-q^{-1}P-(I_m-P),\qquad B=\diag(qI_r,J).
\]
The minus signs leave its positive square unchanged and will produce
the subtraction $UV-P$ needed for integrality. Prescribe the
upper-left block $qI_r$ and off-diagonal blocks $V,U$. Eliminating
these off-diagonal blocks subtracts $U(qI_r)^{-1}V=UV/q$ from the
lower-right block. To leave $J$ as the Schur complement, that block
must therefore be $UV/q+J$. Thus the factorization modulo $q$ is exactly what is
needed to cancel the denominators in this block.

\begin{lemma}\label{lem:integral}
Suppose $P=P^*=P^2\in\Mat_m(\Q[G])$, an odd $D>0$ clears its
denominators, and $\bar P=uv$ factors through $r\ge1$ columns.
Let $q=2^k>1$ satisfy $q\equiv1\pmod D$. Choose arbitrary
integral lifts $U,V$ of the factors supplied by Lemma~\ref{lem:lift}.
Then
\begin{equation}\label{eq:integral-matrix}
 A=\begin{pmatrix}
 qI_r&V\\
 U&q^{-1}\bigl(UV-P-q(I_m-P)\bigr)
 \end{pmatrix}
 \in\Mat_{r+m}(\Z[G])
\end{equation}
has an inverse in $\Mat_{r+m}(\Q[G])$.
\end{lemma}

\begin{proof}
The matrix $P+q(I_m-P)=qI_m+(1-q)P$ is integral because
$D\mid(q-1)$. Its reduction modulo $q$ is $e=P_q$, which is
also the reduction of $UV$. Thus every coefficient of
$UV-P-q(I_m-P)$ is an integer divisible by $q$; the bottom-right
block of Equation~\eqref{eq:integral-matrix} is integral.

For $J$ and $B$ defined above, the complementary projection identities give
$J^{-1}=-qP-(I_m-P)$. For $x\in[0,1]$ put
\[
 L_x=\begin{pmatrix}I_r&0\\xU/q&I_m\end{pmatrix},\qquad
 R_x=\begin{pmatrix}I_r&xV/q\\0&I_m\end{pmatrix}.
\]
Their inverses negate the off-diagonal blocks, and multiplication
without interchanging any factors gives
\begin{equation}\label{eq:path}
 M(x):=L_xBR_x=
 \begin{pmatrix}qI_r&xV\\xU&x^2UV/q+J\end{pmatrix}.
\end{equation}
In particular $A=M(1)$, with the rational group-ring inverse
$R_1^{-1}B^{-1}L_1^{-1}$. For general real $x$, both $M(x)$ and
its inverse lie in the real group ring and depend continuously
differentiably on $x$. No identities $PU=U$, $VP=V$, or
$V=U^*$ are required of the chosen integer lifts.
\end{proof}

The construction is integral at its endpoint, though its simplifying
path need not be integral. We
now show that this path preserves the logarithmic determinant.

\section{The trace calculation}
\label{sec:analysis}

Let $\mathcal A_n$ be the operator-norm closure of
$\Mat_n(\C[G])$ in its left regular representation. For finite
group-ring matrices define
\[
 \tau_n(C)=\sum_{i=1}^n[1]C_{ii}.
\]
The identity-coefficient formula gives
$\tau_n(CD)=\tau_n(DC)$ by summing
$(C_{ij})_g(D_{ji})_{g^{-1}}$ over $i,j,g$. Moreover
$|\tau_n(C)|\le n\|C\|$ and
$\tau_n(C^*)=\overline{\tau_n(C)}$. Thus it extends to a
continuous trace on the unital $C^*$-algebra $\mathcal A_n$.

\begin{lemma}\label{lem:logderivative}
For a continuously differentiable path
$M:[0,1]\to\mathcal A_n$ of invertible elements,
\[
 \frac{d}{dx}\tau_n\bigl(\log(M(x)^*M(x))\bigr)
 =2\operatorname{Re}\tau_n\bigl(M'(x)M(x)^{-1}\bigr).
\]
\end{lemma}

\begin{proof}
Put $N=M^*M$. Continuity of $M$ and its inverse on the compact
interval gives constants $0<\delta<c$ with
$\delta I\le N(x)\le cI$ for all $x$. Write
$W=I-N/c$, so $\|W\|\le\rho=1-\delta/c<1$. Functional calculus
gives the uniformly norm-convergent series
\[
 \log N=(\log c)I-\sum_{j\ge1}\frac{W^j}{j}.
\]
Its derivative series is also uniformly norm-convergent, since
\[
 \left\|\frac{d}{dx}\frac{W^j}{j}\right\|
 \le\sup_{y\in[0,1]}\|W'(y)\|\rho^{j-1}.
\]
In the derivative of $W^j$, cyclicity makes the traces of its
$j$ terms equal; commutation of $N$ and $N'$ is not assumed.
Using $W'=-N'/c$ therefore gives
\[
 \frac{d}{dx}\tau_n(\log N)
 =\frac1c\sum_{j\ge1}\tau_n(W^{j-1}N')
 =\tau_n(N^{-1}N').
\]
Since $N'=M'^*M+M^*M'$, the second summand contributes
$\tau_n(M'M^{-1})$. Cyclicity and
$MN^{-1}=(M^*)^{-1}$ turn the first into
\[
 \tau_n((M^*)^{-1}M'^*)
 =\overline{\tau_n(M'M^{-1})}.
\]
Their sum is the asserted real part. This is the tracial
logarithmic derivative of~\cite[Lemma 2]{FK1952}, written with
the unnormalized trace used here.
\end{proof}

\begin{lemma}\label{lem:shears}
Let $r,m\ge1$ and let $B\in\mathcal A_{r+m}$ be invertible.
Let $S,T$ have respective sizes $m\times r$ and $r\times m$ over
the operator-norm closure of $\C[G]$. For $x\in[0,1]$, let
$L_x,R_x$ be lower and upper block shears with diagonal blocks
$I_r,I_m$ and off-diagonal blocks $xS,xT$, respectively.
Then $\tau_{r+m}(\log((L_xBR_x)^*(L_xBR_x)))$ is constant.
\end{lemma}

\begin{proof}
Write $M=L_xBR_x$. Both shears have inverses obtained by negating
their off-diagonal blocks. The ordered product rule gives
\[
 M'M^{-1}=L_x'L_x^{-1}
 +(L_xB)(R_x'R_x^{-1})(L_xB)^{-1}.
\]
The matrices $L_x'L_x^{-1}$ and $R_x'R_x^{-1}$ have zero diagonal
blocks, so each has trace zero. Cyclicity gives trace zero also for
the conjugate $(L_xB)(R_x'R_x^{-1})(L_xB)^{-1}$.
Lemma~\ref{lem:logderivative} now gives zero derivative.
\end{proof}

\begin{proposition}\label{prop:criterion}
Let $G$ be a discrete group, $m,r\ge1$, and let
$P=P^*=P^2\in\Mat_m(\Q[G])$. Suppose an odd $D>0$ clears its
denominators and $\bar P=uv$ with
$u\in\Mat_{m,r}(\F_2[G])$ and
$v\in\Mat_{r,m}(\F_2[G])$. For every $q=2^k>1$ congruent to one
modulo $D$, the matrix $A$ of Lemma~\ref{lem:integral} satisfies
\begin{align}
 \tau_{r+m}(\log(T_A^*T_A))
   &=2(r-\tau_m(P))\log q,\label{eq:criterion-trace}\\
 \log\det_{\N(G)}(T_A)&=(r-\tau_m(P))\log q.
 \label{eq:criterion-det}
\end{align}
It is integral and rationally invertible, and its positive-square
spectral measure has support in a compact subinterval of $(0,\infty)$.
\end{proposition}

\begin{proof}
Lemma~\ref{lem:integral} has already established integrality and
the finite-support rational inverse. Apply Lemma~\ref{lem:shears}
to Equation~\eqref{eq:path}. At $x=0$ the matrix is
$B=\diag(qI_r,-q^{-1}P-(I_m-P))$. Orthogonality gives
\[
 B^*B=\diag(q^2I_r,q^{-2}P+I_m-P),\qquad
 \log(B^*B)=2\log q\,\diag(I_r,-P).
\]
Taking the trace and using shear invariance proves
Equation~\eqref{eq:criterion-trace}; half this trace is
Equation~\eqref{eq:criterion-det}. The regular representation
of the rational group-ring inverse is bounded, so
\[
 \|T_A^{-1}\|^{-2}I\le T_A^*T_A\le\|T_A\|^2I.
\]
Its logarithm is therefore bounded. In the probability measure
convention from the introduction this also gives
\begin{equation}\label{eq:criterion-integral}
 \int_{(0,\infty)}\log t\,d\mu_A(t)
 =\frac{2(r-\tau_m(P))\log q}{r+m}.
\end{equation}
\end{proof}

\section{Assembly and exact determinant}
\label{sec:assembly}

\begin{proof}[Proof of Theorem~\ref{thm:main}]
Take $K_2,H_2,a_2,b_2$ from Theorem~\ref{thm:input}. Apply
Lemma~\ref{lem:binary} with $d=[K_2:\F_2]\ge1$ and then
Proposition~\ref{prop:separation}. It gives the group
$G_D=(\bigoplus_{H_2}C_3)\rtimes H_2$ and a projection $P$
of size $d+1$ whose binary reduction factors through $d$ columns.
Its trace is $d+\theta$, where
\[
 \theta=\frac13\left(\frac23\right)^{|S|}>0,
 \qquad D=3^{1+|S|}.
\]
The odd integer $D$ clears the coefficients of $P$. Choose
$k\ge1$ with $q=2^k\equiv1\pmod D$.
Proposition~\ref{prop:criterion}, with $r=d$ and $m=d+1$,
produces $A\in\Mat_{2d+1}(\Z[G_D])$ with a rational group-ring
inverse and
\[
 \log\det_{\N(G_D)}(T_A)=-\theta\log q.
\]
Since $q>1$ and $\theta>0$, this proves
\begin{equation}\label{eq:final-det}
 0<\det_{\N(G_D)}(T_A)=q^{-\theta}<1.
\end{equation}
With $n=2d+1$, Equation~\eqref{eq:criterion-integral} reads
\begin{equation}\label{eq:final-integral}
 \int_{(0,\infty)}\log t\,d\mu_A(t)
 =-\frac{2\theta\log q}{2d+1}<0.
\end{equation}
It is finite because of the spectral bounds in
Proposition~\ref{prop:criterion}. Finally, a finite generating
set of $H_2$ together with $z_1$ generates $G_D$, since
$z_h=hz_1h^{-1}$. Thus $G_D$ is finitely generated.
\end{proof}

The target group differs from the companion's group and contains
the cyclic coordinates of order three. The construction has a
finite negative logarithmic determinant and no zero spectral mass;
it is a failure of the asserted lower bound, not of determinant
class. Applying the sofic determinant bound of
\cite[Theorem 5]{ElekSzabo2005} to the positive integral matrix
$A^*A$ proves that $G_D$ is not sofic. This last implication is an
external consequence, not an ingredient in constructing $A$.

\begingroup
\raggedright
\bibliographystyle{plain}
\bibliography{references}
\endgroup
\end{document}
